\subsection{零化子。忠实模。单生成模。}

定义11. A模E的子集S的零化子是所有满足下列条件的元素 \(\alpha \in \mathbf{A}\) 组成的集合：\(\alpha x = 0\) 对所有 \(x \in S\) 成立.

S 的零化子通常记为 Ann(S)；对于由单个元素 x 组成的子集 S，我们写 Ann(x) 而不写 Ann(\{x\})，并把 Ann(x) 称为 x 的零化子。

关系 \(\alpha x = 0\) 也可以说成： \(x\) 被 \(\alpha\) 零化.

立即可以看出，E的任意子集S的零化子是A的一个左理想；要使它等于A，必须且只需（根据 \((\mathbf{M}_{\mathrm{IV}})\) ）有 \(S = \{0\}\) 。如果 E 的两个子集 S、T 满足 \(S \subset T\) ，则 T 的零化子包含于 S 的零化子中。如果 \((\mathbf{S}_{t})_{t \in I}\) 是E的任意子集族，则并集 \(\bigcup_{t} S_{t}\) 的零化子是诸 \(S_{t}\) 的零化子的交。特别地，E 的子集 S 的零化子是 S 中各元素的零化子的交。说 E 的一个元素是自由的，等价于说它的零化子是 \(\{0\}\) 。对所有 \(x \in E\) 以及所有 \(\alpha \in A\) ， \(\alpha x\) 的零化子是这样的 \(\beta \in A\) 组成的集合，使得 \(\beta \alpha \in \text{Ann}(x)\) .

子模M在E中的零化子是A的一个双边理想；因为，若 \(\alpha x = 0\) 对所有 \(x \in M\) 成立，则 \(\alpha(\beta x) = 0\) 也对所有 \(x \in M\) 以及所有 \(\beta \in A\) 成立，因此 \(\alpha\beta\) 对所有 \(\beta \in A\) 都属于M的零化子。特别地，E的零化子是A的一个双边理想。

对于所有 \(\alpha \in \mathbf{A}\) ，设 \(h_{\alpha}\) 为位似变换 \(x \mapsto \alpha x\) ；已知映射 \(\alpha \mapsto h_{\alpha}\) （从 \(\mathbf{A}\) 到自同态环 \(\mathcal{E} = \operatorname{Hom}_{\mathbf{Z}}(\mathbf{E}, \mathbf{E})\) ，即无算子的交换群 \(\mathbf{E}\) 的自同态环）是一个环同态（ \(\S 2\) ，第5号）。该同态下0的原像是零化子 \(\mathfrak{a}\) （即 \(\mathbf{E}\) 的零化子）；因此， \(\mathbf{A}\) 在同态 \(\alpha \mapsto h_{\alpha}\) 下的像同构于商环 \(\mathbf{A}/\mathfrak{a}\) 。模 \(\mathbf{E}\) 称为忠实的，如果其零化子 \(\mathfrak{a}\) 缩减为0。

设 E 为任意 A-模，a 为 A 中包含于 Ann(E) 的一个双边理想，并设 \(\dot{\alpha}\) 为商环 A/a 的一个元素；对所有 \(x \in E\) ，元素 \(\alpha x\) 只依赖于 \(\alpha \in A\) 所属的类 \(\dot{\alpha} \mod a\) ；若将此元素记为 \(\dot{\alpha}x\) ，则立即可以看出，映射 \((\dot{\alpha}, x) \mapsto \dot{\alpha}x\) （连同 E 上的加法）在 E 上定义了一个 (A/a)-模结构。当 \(a = \text{Ann}(E)\) 时，如此定义的 (A/a)-模 E 是忠实的；我们称它为与 A-模 E 相伴的忠实模。注意，A-模 E 的每个子模也是相伴忠实模的子模，反之亦然。

定义12. 一个模称为单生成的，如果它由单个元素生成。

第7号命题9表明，如果E是一个单生成A-模且a是生成E的元素，则E与集合A.a 相同，它由 \(\xi a\) 组成，其中 \(\xi\) 遍历A。

例. (1) 每个单生成群，由于是交换的（I，§ 4，第10号，命题18），都是一个单生成Z-模。

\begin{enumerate}
\def\labelenumi{(\arabic{enumi})}
\setcounter{enumi}{1}
\item
  若A是交换环，则A-模A的单生成子模恰好就是环A的主理想（I，§ 8，第6号）。
\item
  每个单A-模E是单生成的，因为由E中元素 \(\neq0\) 生成的E的子模必然等于E。
\end{enumerate}

命题22. 设 A 为一个环。 \(\mathbf{A}_s\) 的每个商模都是单生成的。反之，设 E 为一个单生成 A-模， \(c\) 为 E 的一个生成元，a 为其零化子；线性映射 \(\xi \mapsto \xi c\) 在过渡到商后定义了 \(\mathbf{A}_s / \mathfrak{a}\) 到 E 上的一个同构。

由于 \(A_{s}\) 本身是单生成的（由 1 生成），第一个断言由第7号命题9的推论1得出。第二个断言是显然的，因为 \(\xi \mapsto \xi c\) 根据假设是满射的且核为 a。

注意，若 A 非交换，则单生成 A-模 E 的两个不同生成元 c, \(c'\) 的零化子一般互不相同，并且也与模 E 的零化子不同。另一方面，若 A 是交换的，则 E 的一个生成元 c 的零化子包含于 E 的每个元素的零化子之中，因此就是整个 E 的零化子。

推论。单生成A-模E的每个子模都同构于一个商模b/a，其中a和b是A的两个左理想，使得 \(a \subset b\) 。单生成A-模的每个商模都是单生成的。

第二个断言是显然的，第一个断言由命题22及I，§4，第6号，定理4得出。

另一方面，请注意，单生成模的子模不一定是单生成的。例如，若A是一个存在非主理想的交换环（VII，§1，第1号），则这些理想是单生成A-模A的非单生成子模。

由定义可知，A-模E的由E中元素族 \((a_{t})\) 生成的子模是诸单生成子模

\(Aa_{t}\) 之和；要使 \((a_{t})\) 成为E的一组基，必须且只需每一个 \(a_{t}\) 都是E的自由元素，并且诸 \(Aa_{t}\) 的和是直和。

命题23. 设E是一个A-模，它是一族无限多个非零子模 \((\mathbf{M}_t)_{t\in I}\) 的直和。对于E的每一个生成系S， \(\operatorname{Card}(\mathbf{S})\geqslant \operatorname{Card}(\mathbf{I})\)

对于所有 \(x \in S\) ，设 \(C_x\) 为指标 \(i \in I\) 组成的有限集，使得 \(x\) 在 \(M_i\) 中的分量 \(\neq 0\) ，并设 \(C = \bigcup_{x \in S} C_x\) 。每个 \(x \in S\) 按定义属于 \(E\) 的由诸 \(M_i\) （ \(i \in C\) ）组成的直和子模，且 \(S\) 生成 \(E\) 这一假设因此蕴含 \(C = I\) ；由于 \(I\) 按假设是无限的， \(S\) 也是无限的（集合论，III，§5，第1号，命题1的推论1）；因此 \(\text{Card}(I) = \text{Card}(C) \leqslant \text{Card}(S)\) （集合论，III，§6，第3号，定理2的推论3）。

推论1. 在命题23的假设下，再设每个 \(\mathbf{M}_{\iota}\) 都是单生成的，且 \(\mathbf{E}\) 是第二个由非零单生成子模组成的族 \((\mathrm{N}_{\lambda})_{\lambda \in \mathbf{L}}\) 的直和。则 \(\operatorname{Card}(\mathbf{L}) = \operatorname{Card}(\mathbf{I})\) .

若 \(b_{\lambda}\) 是 \(\mathbf{N}_{\lambda}\) 的一个生成元，则诸 \(b_{\lambda}\) 组成的集合是 \(\mathbf{E}\) 的一个生成系，因此 \(\operatorname{Card}(\mathbf{L}) \geqslant \operatorname{Card}(\mathbf{I})\) 。特别地， \(\mathbf{L}\) 是无限的，并且交换 \((\mathbf{M}_i)\) 与 \((\mathbf{N}_{\lambda})\) 的角色，类似地有 \(\operatorname{Card}(\mathbf{I}) \geqslant \operatorname{Card}(\mathbf{L})\) ，由此得出推论。

推论2. 若模E具有无限基B，则E的每个生成系的基数 \(\geqslant\) Card(B)，且E的每个基都与B等势。

